\documentclass[12pt,a4paper]{article}
\usepackage[utf8]{inputenc}
\usepackage[T1]{fontenc}
\usepackage{geometry}
\usepackage{setspace}
\usepackage{graphicx}
\usepackage{booktabs}
\usepackage{float}
\usepackage{amsmath}
\usepackage{tabularx}
\usepackage{longtable}
\usepackage{array}
\geometry{margin=1in}
\onehalfspacing
\setcounter{tocdepth}{2}
\renewcommand{\arraystretch}{1.15}

\title{Natural Language Processing Project Report\\Sentiment Analysis for E-commerce Reviews using Machine Learning Techniques.}
\author{Team SkyDom+}
\date{February 2026}

\begin{document}

\begin{titlepage}
\centering
\vspace*{1.0cm}

\includegraphics[width=0.62\textwidth]{p01_chart1.png}\par
\vspace{1.0cm}
{\LARGE Sentiment Analysis for E-commerce Reviews\par}
\vspace{0.2cm}
{\LARGE using Machine Learning Techniques.\par}
\vspace{0.8cm}
{\Large Case Study Report\par}
\vspace{0.8cm}
{\large Team: Team SkyDom+\par}
\vspace{0.3cm}
{\large Group members:\par}
\vspace{0.4cm}
{\large
\begin{tabular}{@{}>{\raggedright\arraybackslash}p{5.8cm} >{\raggedright\arraybackslash}p{2.7cm}@{}}
Nguyen Nhu Tien Tuan & 22BA13314 \\
Tran Hien Chuong & 22BA13056 \\
Do Gia Bao & 22BA13042 \\
Bui The Anh & 22BA13007 \\
Le Duc An & 22BA13002 \\
Pham Quang Tung & 22BA13320 \\
Khuong Bao Duc & 22BA13077 \\
\end{tabular}\par}
\vfill
{\large February 25, 2026\par}
\end{titlepage}

\pagenumbering{roman}

\begin{abstract}
This report presents an NLP project for review understanding on Amazon Gift Card data. The project has three NLP layers. The first layer is sentence-level sentiment classification with context-aware text features. The second layer is multi-label issue extraction to identify concrete complaint types. The third layer is a transformer extension for deep contextual semantics and hard-case analysis. The work uses 152,410 raw reviews, with strong sentiment labels from ratings 1, 2, 4, and 5, and an uncertain analysis set from rating 3. The selected classic sentiment variant uses word and character TF-IDF with negation tagging. On test data, it reaches precision for negative class of 0.677, recall for negative class of 0.932, and F2 for negative class of 0.867. The transformer extension reaches higher overall accuracy and higher negative precision, but lower raw negative recall at fixed threshold. For issue extraction, the word-level multi-label model gets test micro-F1 of 0.765, and the word plus character model gets 0.763. To improve course alignment without adding data, the project now includes a strong syllabus upgrade package with Naive Bayes, Perceptron, LinearSVM, SGD log-loss, semantic SVD baselines, and a feed-forward MLP baseline. In that package, test recall\_0 ranges from 0.719 to 0.941, with the best F2\_0 at 0.855 for balanced Logistic Regression and the best recall\_0 at 0.941 for Complement Naive Bayes. The package also adds unigram and bigram language models with add-k smoothing, where bigram test perplexity improves from 377.45 to 362.72. A course-fit matrix built from real artifacts reports estimated syllabus coverage of 75.7\%.
\end{abstract}

\tableofcontents
\newpage
\pagenumbering{arabic}

\section{Introduction}
Natural language in user reviews is noisy, short, and mixed. A single review can include praise, complaint, slang, typo, and domain words in the same sentence. This makes review understanding a good NLP problem for both classic and modern methods \cite{jurafsky2026,turney2002}.

This project studies review text from Amazon Gift Cards. The goal is not only to assign a positive or negative label. The goal is to build a practical language pipeline that can understand sentiment, mark uncertain cases, and extract concrete issue tags for support work.

The project is designed as an end-to-end NLP system. It includes text normalization, lexical and subword feature modeling, context-aware handling of negation and contrast, supervised learning, uncertainty control, multi-label issue extraction, and transformer-based extension.

\section{NLP Problem Definition}
\subsection{Task A: Sentiment Understanding}
Task A is binary sentiment classification from review text. Ratings 1 and 2 are mapped to negative class. Ratings 4 and 5 are mapped to positive class. Ratings with value 3 are used for uncertainty analysis and are not part of strong supervised training \cite{turney2002,wang2012}.

The practical decision policy is not plain accuracy-first. In support triage, missing a bad review is risky. Therefore model selection follows an internal negative-first heuristic: prioritize recall for class 0, then F2 for class 0, then precision for class 0.

\subsection{Task B: Issue Extraction}
Task B is multi-label issue extraction. One review can contain more than one issue, so single-label tagging is not enough \cite{tsoumakas2011,zhang2014}. The output label set in this project has nine issue types and allows multi-label assignment per row.

This task is a core NLP mapping from unstructured review text to structured issue signals. It helps support teams move from general sentiment to action categories.

\subsection{Task C: Deep Context Extension}
Task C is transformer fine-tuning as NLP extension. The extension aims to test deep contextual language modeling for difficult short phrases and negation patterns. It also provides a direct comparison with classic sparse representations \cite{vaswani2017,devlin2019}.

\section{Dataset and Language Characteristics}
\subsection{Data Source and Scale}
The source file contains 152,410 reviews with text and rating fields. Table \ref{tab:nlp_rating_distribution} shows full rating distribution.

\begin{table}[H]
\centering
\caption{Rating distribution in source data}
\label{tab:nlp_rating_distribution}
\begin{tabular}{lr}
\toprule
Rating & Number of reviews \\
\midrule
1 & 12,326 \\
2 & 1,873 \\
3 & 3,271 \\
4 & 6,692 \\
5 & 128,248 \\
\bottomrule
\end{tabular}
\end{table}

Strong sentiment labels include ratings 1, 2, 4, and 5. This gives 149,139 rows. Negative class has 14,199 rows and positive class has 134,940 rows. The class ratio is about 9.5 to 1. This imbalance is important for NLP modeling choices and for precision-recall focused evaluation \cite{saito2015}.

\subsection{Observed Language Noise}
The review language contains common NLP noise forms. The dataset includes short forms such as gr8, thx, idk, and informal phrases. It includes mixed polarity in one sentence, such as good but late delivery. It includes clear negation phrases such as not bad and not good.

The dataset also has domain-specific words for gift cards, code redemption, activation, delivery, support, and refund. These patterns motivate both lexical mapping and data-driven learning.

\subsection{Issue Label Data}
For issue extraction, the manually labeled dataset has 19,000 rows. Table \ref{tab:nlp_issue_frequency} shows label frequency.

\begin{table}[H]
\centering
\caption{Issue label frequency in manual NLP annotation set}
\label{tab:nlp_issue_frequency}
\begin{tabular}{lrr}
\toprule
Label & Count & Share \\
\midrule
delivery\_shipping & 972 & 0.0512 \\
redemption\_activation & 284 & 0.0149 \\
product\_quality & 1,251 & 0.0658 \\
customer\_service & 182 & 0.0096 \\
refund\_return & 93 & 0.0049 \\
usability & 873 & 0.0459 \\
value\_price & 558 & 0.0294 \\
fraud\_scam & 144 & 0.0076 \\
other & 15,960 & 0.8400 \\
\bottomrule
\end{tabular}
\end{table}

Average labels per review is 1.069. This supports multi-label modeling because many rows are not single-concept text.

\section{Text Processing Design}
\subsection{Normalization}
The core normalization pipeline uses lowercase conversion, URL removal, punctuation cleanup, and white-space normalization. This reduces token sparsity while keeping core sentiment words.

Optional abbreviation normalization maps common short forms to full forms. This includes mappings such as gr8 to great, thx to thanks, u to you, and idk to i do not know.

\subsection{Negation Handling}
Negation is a key NLP challenge in short reviews. The project uses a negation tagging method. After a negation cue, the next few tokens receive a NOT\_ prefix. The window size is three tokens by default \cite{wiegand2010}.

This method is simple but useful for sparse models. It helps separate good from NOT\_good and bad from NOT\_bad in feature space.

\subsection{Contrast Handling}
The pipeline also models simple discourse contrast. When markers such as but, however, although, though, and yet appear, text can be split into left and right clauses. Separate clause features and marker flags are then added.

This design captures a common review pattern where the second clause carries the final user opinion.

\subsection{Character-level Modeling}
Character 3 to 5 gram features are used as optional union with word features. This helps with typo, short slang, and small spelling variants. In NLP tasks on user text, this often improves robustness \cite{joulin2017}.

\section{Feature Engineering for NLP}
\subsection{Classic Sparse Representation}
The main representation is TF-IDF with word n-grams. The baseline uses 1 to 2 grams. Other variants test 1 to 3 grams, negation tags, contrast clauses, character grams, and lexicon features \cite{manning2008ir,wang2012}.

Table \ref{tab:nlp_variant_design} summarizes tested context variants.

\begin{table}[H]
\centering
\caption{NLP context feature variants}
\label{tab:nlp_variant_design}
\begin{tabular}{>{\raggedright\arraybackslash}p{1.2cm} >{\raggedright\arraybackslash}p{12cm}}
\toprule
Variant & Feature design \\
\midrule
V0 & Baseline word TF-IDF with n-gram range 1 to 2 \\
V1 & Word TF-IDF with n-gram range 1 to 3 \\
V2 & Word 1 to 2 plus negation tagging \\
V3 & Word 1 to 2 plus clause split with contrast marker \\
V4 & Word 1 to 2 plus negation tagging plus clause split \\
V5 & Word 1 to 2 plus character 3 to 5 grams \\
V6 & Word 1 to 2 plus character 3 to 5 grams plus negation tagging \\
V7 & Word 1 to 2 plus negation tagging plus sentiment lexicon block \\
\bottomrule
\end{tabular}
\end{table}

\subsection{Lexicon Feature Block}
The project includes an AFINN-based sentiment lexicon block in one variant. It computes raw score, negation-aware score, intensity-aware score, and contrast-weighted score. This adds explicit lexical prior to sparse representation \cite{nielsen2011}.

This part is useful for NLP course learning because it connects classic lexicon methods with modern supervised learning in one pipeline.

\section{Modeling Approach}
\subsection{Sentiment Classifiers}
The main sentiment classifier is Logistic Regression. Chi-square feature selection is applied with top K values from 1,000 to 10,000. Class weight settings include none, balanced, and weighted variants \cite{forman2003}.

Additional sentiment models include L1 Logistic Regression for sparse embedded behavior, Decision Tree, and Random Forest.

\subsection{Issue Extraction Model}
Issue extraction uses one-vs-rest training across nine labels. The default model is Logistic Regression with balanced class weight. A LinearSVM baseline is also available \cite{rifkin2004}.

Threshold tuning is done per label on validation data. This is important in multi-label NLP because one global threshold is often not enough for all labels.

\subsection{Transformer Extension}
The transformer extension fine-tunes a sequence classifier with class-weighted loss. It uses cleaned text and evaluates with the same uncertainty and threshold framework for direct comparison \cite{devlin2019}.

This extension gives deeper contextual semantics and helps analyze hard short phrases where sparse models may fail.

\section{Experimental Setup}
\subsection{Split and Protocol}
All experiments use seed 42. Strong sentiment labels are split with train, validation, and test ratio 70, 10, and 20. Three-star reviews are held for uncertainty analysis.

Issue extraction uses 19,000 labeled rows with split sizes 13,300 for train, 1,900 for validation, and 3,800 for test in saved runs.

\subsection{Main Metrics}
For sentiment, key metrics are precision\_0, recall\_0, F1\_0, and F2\_0, along with full accuracy and F1. Because the strong-label set is imbalanced, model trade-offs are interpreted mainly through precision-recall behavior rather than ROC-only views \cite{saito2015}. Selective metrics under uncertainty include coverage, uncertain rate, selective precision\_0, selective recall\_0, selective F2\_0, and false-negative rate on covered samples.

For issue extraction, key metrics are micro-F1, macro-F1, subset accuracy, hamming loss, and per-label precision, recall, and F1.

\subsection{Uncertainty Rule}
The uncertainty policy uses low threshold and high threshold on positive probability. It follows a selective classification abstention setup and can return uncertain in the middle range \cite{geifman2017}. We apply an additional heuristic in the classic pipeline that returns uncertain for too short or sparse text cases.

This rule makes NLP output safer for production because low-information text is not forced into a hard wrong class.

\section{Team Workflow and Quality Assurance}
\subsection{Module Decomposition and Responsibility}
To keep the implementation scalable for a seven-member team, the project is split into clear technical modules with explicit ownership and cross-check requirements. Table \ref{tab:nlp_team_modules} summarizes the structure used during development and defense preparation.

\begin{table}[H]
\centering
\caption{Team module split and quality responsibility}
\label{tab:nlp_team_modules}
\begin{tabular}{>{\raggedright\arraybackslash}p{3.2cm} >{\raggedright\arraybackslash}p{4.8cm} >{\raggedright\arraybackslash}p{3.8cm} >{\raggedright\arraybackslash}p{2.6cm}}
\toprule
Module & Main responsibility & Cross-check by non-owner & Main artifact path \\
\midrule
Sentiment core & cleaning, TF-IDF features, Chi-square, logistic and uncertainty policy & rerun from raw JSONL and verify metric consistency & \texttt{results/dm2\_steps/} \\
Issue extraction & annotation schema validation, OVR training, per-label thresholds & verify label distribution and per-label F1 deltas & \texttt{results/issue\_steps/} \\
Transformer extension & fine-tuning, threshold sweep, hard-case comparison & compare against classic using same evaluation protocol & \texttt{results/nlp\_ext/} \\
Reporting and demo & report writing, figure integration, demo scenario rehearsal & final end-to-end rerun and Q\&A preparation & \texttt{results/reports/} \\
\bottomrule
\end{tabular}
\end{table}

\subsection{Cross-understanding, Review, and Reproducibility}
The workflow is designed so that each critical module can be executed by at least one member who is not the original author. For each module handover, the receiving member reruns the pipeline from saved commands and checks three gates: no train-test leakage, fixed random seed behavior, and consistency between table values and exported CSV files.

For code quality control, we use a branch-and-review policy with one non-author review per substantial change. The review checklist focuses on metric correctness, split integrity, and reproducibility. This process aligns with course criteria on module split, shared understanding, and cross-review practice.

\subsection{Defense Evidence Bundle}
To make teamwork evidence auditable in the defense session, the bundle includes: role/module matrix, command logs for reruns, model-output CSV artifacts, and review notes linked to changed modules. This evidence complements the quantitative model results and helps justify process quality, not only final scores.

\section{No-New-Data Syllabus Upgrade Package}
The project now includes an extra package that improves alignment with NLP course topics without adding any new labeled data. The package reuses the same train, validation, and test split, and keeps the same text cleaning policy.

The first upgrade block is a strong classic benchmark suite. It includes balanced Logistic Regression, Multinomial Naive Bayes, Complement Naive Bayes, Perceptron, LinearSVM, SGD with log-loss, semantic Logistic Regression on SVD vectors, and a feed-forward MLP on the same SVD vectors. The train size is capped at 35,000 rows for this block to keep runtime practical while preserving class ratio.

The second upgrade block is a language model block. It trains unigram and bigram language models with add-k smoothing and reports perplexity on validation and test splits. This gives direct evidence for the N-gram language model topic in the course \cite{jurafsky2026,manning2008ir}.

The third upgrade block builds a topic-coverage matrix from project artifacts. The matrix uses explicit topic scores and reports an overall syllabus coverage score for defense discussion.

\section{Results for Classic NLP Sentiment Model}
\subsection{Baseline Selection by K}
The baseline sweep on V0 shows best settings per K under negative-first rule. Table \ref{tab:nlp_chi2_best_per_k} reports selected rows from validation.

\begin{table}[H]
\centering
\caption{Best baseline sentiment result per K on validation}
\label{tab:nlp_chi2_best_per_k}
\begin{tabular}{lcccccc}
\toprule
K & Class weight & Precision\_0 & Recall\_0 & F2\_0 & Accuracy & F1 \\
\midrule
1,000 & w10 & 0.555 & 0.910 & 0.806 & 0.922 & 0.955 \\
2,000 & w10 & 0.584 & 0.920 & 0.825 & 0.930 & 0.960 \\
5,000 & balanced & 0.606 & 0.935 & 0.843 & 0.936 & 0.963 \\
10,000 & balanced & 0.621 & 0.936 & 0.850 & 0.940 & 0.966 \\
\bottomrule
\end{tabular}
\end{table}

\subsection{Variant Comparison}
After full variant sweep, selected classic setting is V6 with K equal to 10,000 and class weight w10. Table \ref{tab:nlp_variant_test} shows test comparison across variants.

\begin{table}[H]
\centering
\caption{Classic NLP sentiment variants on test data}
\label{tab:nlp_variant_test}
\begin{tabular}{lcccccc}
\toprule
Variant & K & Weight & Precision\_0 & Recall\_0 & F2\_0 & Coverage \\
\midrule
V0 & 10,000 & balanced & 0.638 & 0.937 & 0.856 & 0.870 \\
V1 & 10,000 & balanced & 0.643 & 0.932 & 0.855 & 0.871 \\
V2 & 5,000 & w10 & 0.647 & 0.926 & 0.852 & 0.874 \\
V3 & 10,000 & balanced & 0.626 & 0.923 & 0.843 & 0.868 \\
V4 & 10,000 & balanced & 0.635 & 0.921 & 0.845 & 0.872 \\
V5 & 10,000 & balanced & 0.653 & 0.933 & 0.859 & 0.878 \\
V6 & 10,000 & w10 & 0.677 & 0.932 & 0.867 & 0.885 \\
V7 & 5,000 & w10 & 0.657 & 0.927 & 0.857 & 0.880 \\
\bottomrule
\end{tabular}
\end{table}

V6 offers a good balance for NLP sentiment in this project. It keeps high negative recall and improves negative precision compared with many other variants.

\subsection{Threshold and Uncertainty Behavior}
For selected classic model V6, default threshold 0.40 and 0.60 gives coverage 0.885 and selective recall\_0 0.956. Table \ref{tab:nlp_classic_threshold} shows the threshold sweep.

\begin{table}[H]
\centering
\caption{Threshold sweep for selected classic NLP sentiment model}
\label{tab:nlp_classic_threshold}
\begin{tabular}{lcccc}
\toprule
Threshold pair & Coverage & Selective precision\_0 & Selective recall\_0 & Selective F2\_0 \\
\midrule
0.40 / 0.60 & 0.885 & 0.730 & 0.956 & 0.900 \\
0.45 / 0.60 & 0.891 & 0.708 & 0.956 & 0.894 \\
0.50 / 0.65 & 0.886 & 0.681 & 0.963 & 0.889 \\
0.55 / 0.70 & 0.881 & 0.646 & 0.970 & 0.882 \\
\bottomrule
\end{tabular}
\end{table}

\section{Results for Transformer NLP Extension}
\subsection{Transformer Metrics}
Transformer fine-tuning results are saved in NLP extension outputs. Table \ref{tab:nlp_transformer_metrics} summarizes validation and test values.

\begin{table}[H]
\centering
\caption{Transformer NLP extension metrics}
\label{tab:nlp_transformer_metrics}
\begin{tabular}{lcccccc}
\toprule
Split & Accuracy & F1 & Precision\_0 & Recall\_0 & Coverage & Selective recall\_0 \\
\midrule
Validation & 0.968 & 0.982 & 0.821 & 0.854 & 0.908 & 0.862 \\
Test & 0.967 & 0.982 & 0.822 & 0.833 & 0.910 & 0.846 \\
\bottomrule
\end{tabular}
\end{table}

Compared with classic V6, transformer gives higher overall accuracy and higher negative precision at base threshold. Classic V6 still keeps much higher raw negative recall.

\subsection{Transformer Threshold Sweep}
Table \ref{tab:nlp_transformer_threshold} shows transformer selective behavior under threshold sweep.

\begin{table}[H]
\centering
\caption{Transformer threshold sweep on test data}
\label{tab:nlp_transformer_threshold}
\begin{tabular}{lcccc}
\toprule
Threshold pair & Coverage & Selective precision\_0 & Selective recall\_0 & Selective F2\_0 \\
\midrule
0.40 / 0.60 & 0.910 & 0.831 & 0.846 & 0.843 \\
0.45 / 0.60 & 0.911 & 0.827 & 0.847 & 0.843 \\
0.50 / 0.65 & 0.912 & 0.821 & 0.850 & 0.844 \\
0.55 / 0.70 & 0.911 & 0.818 & 0.855 & 0.847 \\
\bottomrule
\end{tabular}
\end{table}

The transformer curve is stable, with coverage near 0.91 and moderate gain in selective recall when threshold pair moves to 0.55 and 0.70.

\subsection{Hard-case Comparison}
Hard-case results in Table \ref{tab:nlp_hard_cases} compare classic baseline and transformer decisions.

\begin{table}[H]
\centering
\caption{Hard-case comparison between classic and transformer sentiment models}
\label{tab:nlp_hard_cases}
\begin{tabular}{>{\raggedright\arraybackslash}p{4.1cm} >{\centering\arraybackslash}p{2.3cm} >{\centering\arraybackslash}p{2.3cm} >{\raggedright\arraybackslash}p{4.3cm}}
\toprule
Text & Baseline decision & Transformer decision & Main observation \\
\midrule
not bad & Negative & Positive & transformer reads idiom as positive \\
not good & Negative & Negative & both models agree on negative \\
good but late delivery & Negative & Positive & transformer weights first positive clause more \\
gr8 & Uncertain & Uncertain & both return uncertain short text \\
thx & Uncertain & Uncertain & both return uncertain short text \\
idk & Uncertain & Uncertain & both return uncertain short text \\
\bottomrule
\end{tabular}
\end{table}

This table shows that transformer handles some semantic patterns differently. It can improve idiomatic cases such as not bad, but it can also move mixed complaint text toward positive.

\section{Results for NLP Issue Extraction}
\subsection{Overall Multi-label Results}
Table \ref{tab:nlp_issue_overall} compares three issue extraction setups.

\begin{table}[H]
\centering
\small
\setlength{\tabcolsep}{4pt}
\caption{NLP issue extraction overall metrics}
\label{tab:nlp_issue_overall}
\begin{tabularx}{\textwidth}{
>{\raggedright\arraybackslash}X
>{\centering\arraybackslash}p{2.1cm}
>{\centering\arraybackslash}p{1.8cm}
>{\centering\arraybackslash}p{1.8cm}
>{\centering\arraybackslash}p{2.4cm}}
\toprule
Setup & Val micro-F1 & Test micro-F1 & Test macro-F1 & Test subset accuracy \\
\midrule
Word TF-IDF plus OVR Logistic & 0.759 & 0.765 & 0.381 & 0.606 \\
Word TF-IDF plus OVR LinearSVM & 0.750 & 0.757 & 0.371 & 0.602 \\
Word plus Char TF-IDF plus OVR Logistic & 0.764 & 0.763 & 0.408 & 0.603 \\
\bottomrule
\end{tabularx}
\end{table}

Word plus character representation gives higher macro-F1, while the word-only setup keeps slightly higher micro-F1 and subset accuracy on this run.

\subsection{Per-label Improvement}
Table \ref{tab:nlp_issue_per_label_delta} shows per-label test F1 change from word to word plus char setup.

\begin{table}[H]
\centering
\caption{Per-label test F1 change in issue extraction}
\label{tab:nlp_issue_per_label_delta}
\begin{tabular}{lccc}
\toprule
Label & F1 word & F1 word plus char & Delta \\
\midrule
delivery\_shipping & 0.475 & 0.462 & -0.013 \\
redemption\_activation & 0.267 & 0.304 & +0.037 \\
product\_quality & 0.308 & 0.301 & -0.007 \\
customer\_service & 0.314 & 0.185 & -0.128 \\
refund\_return & 0.145 & 0.333 & +0.188 \\
usability & 0.462 & 0.476 & +0.014 \\
value\_price & 0.284 & 0.316 & +0.032 \\
fraud\_scam & 0.265 & 0.388 & +0.123 \\
other & 0.911 & 0.910 & -0.001 \\
\bottomrule
\end{tabular}
\end{table}

Character features help some sparse labels such as refund\_return and fraud\_scam, but can reduce F1 on a few labels. This indicates a precision-recall trade-off by label.

\subsection{Threshold Strategy in Multi-label Setting}
Issue extraction uses per-label tuned thresholds. Threshold profiles differ between word-only and word+char setups because score distributions change by label.

This behavior is expected in NLP multi-label tasks where each label has different prior frequency and text signal strength.

\section{Results for Syllabus Upgrade Package}
\subsection{Classic Baselines for Core Course Algorithms}
Table \ref{tab:nlp_syllabus_bench} reports test metrics for the added baseline block. This block uses the same data split and does not add new labels.

\begin{table}[H]
\centering
\caption{Test results for added NLP syllabus baselines}
\label{tab:nlp_syllabus_bench}
\begin{tabular}{lccc}
\toprule
Model & Precision\_0 & Recall\_0 & F2\_0 \\
\midrule
Logistic Regression (balanced) & 0.699 & 0.906 & 0.855 \\
Multinomial Naive Bayes & 0.663 & 0.895 & 0.836 \\
Complement Naive Bayes & 0.531 & 0.941 & 0.815 \\
Perceptron & 0.686 & 0.842 & 0.806 \\
LinearSVM & 0.795 & 0.829 & 0.822 \\
SGD log-loss & 0.618 & 0.915 & 0.835 \\
SVD semantic Logistic Regression & 0.505 & 0.911 & 0.785 \\
FFNN MLP on SVD vectors & 0.796 & 0.719 & 0.733 \\
\bottomrule
\end{tabular}
\end{table}

These results show clear behavior differences. Complement Naive Bayes has the highest negative recall at 0.941, but precision is low at 0.531. Balanced Logistic Regression gives the best F2\_0 at 0.855 with stronger precision. LinearSVM has the highest precision\_0 at 0.795 but lower recall\_0 than logistic. This is useful for exam and defense because it shows the trade-off between models from different course topics under one dataset.

\subsection{Selective Behavior in the Added Baselines}
For probability models in the added block, the same uncertainty band 0.40 and 0.60 is applied. Balanced Logistic Regression gives test coverage 0.883 and selective recall\_0 0.932. Complement Naive Bayes gives test coverage 0.900 and selective recall\_0 0.952, but selective precision\_0 is only 0.568. This confirms that the new block is consistent with the project policy and still keeps transparent trade-off control.

\subsection{N-Gram Language Model Results}
Table \ref{tab:nlp_ngram_results} reports perplexity for unigram and bigram models with add-k equal to 1.0.

\begin{table}[H]
\centering
\caption{Perplexity for N-gram language model block}
\label{tab:nlp_ngram_results}
\begin{tabular}{lcc}
\toprule
Model & Validation perplexity & Test perplexity \\
\midrule
Unigram add-k & 375.91 & 377.45 \\
Bigram add-k & 358.24 & 362.72 \\
\bottomrule
\end{tabular}
\end{table}

Bigram gives lower perplexity than unigram on both splits. This is expected because bigram uses local context from the previous token. The gap is not large because review text is short and noisy, but it is stable and meaningful for the language model topic.

\subsection{Updated Course-Fit Matrix}
Table \ref{tab:nlp_course_fit} shows the updated topic-coverage matrix after the new upgrade block.

\begin{table}[H]
\centering
\caption{Updated topic coverage after no-new-data upgrade}
\label{tab:nlp_course_fit}
\begin{tabular}{lcc}
\toprule
Topic group & Coverage percent & Status \\
\midrule
Probability and evaluation & 100 & strong \\
N-gram language models & 100 & strong \\
Naive Bayes and Perceptron & 100 & strong \\
Logistic and linear models & 100 & strong \\
Feed-forward neural network & 100 & strong \\
Vector semantics and embeddings & 80 & medium-strong \\
Transformers and LLM basics & 100 & strong \\
Pytorch practice & 60 & medium \\
RNN and LSTM & 20 & weak \\
Masked language models & 20 & weak \\
LLM applications & 30 & weak \\
\bottomrule
\end{tabular}
\end{table}

The estimated overall coverage score is 75.7\%. This is much higher than the previous state where core classification topics were strong but many exam topics were still missing from the artifact set.

\section{Classic Versus Transformer Discussion}
Classic sparse NLP and transformer NLP provide different strengths. Classic V6 keeps high negative recall and strong selective recall in this project setting. Transformer gives higher global accuracy and stronger negative precision.

For support triage that needs to catch as many negatives as possible, classic V6 is more aligned with the target policy. For semantic nuance and idiom handling, transformer can be useful and can be combined with rule-based control.

A practical hybrid path is possible. One option is to keep classic model as primary triage model and use transformer as secondary scorer on hard uncertain cases.

\section{Error Analysis and Failure Modes}
Table \ref{tab:nlp_error_cases} gives concrete wrong or risky prediction cases with model-specific explanation. This addresses error analysis beyond aggregate metrics.

\begin{table}[H]
\centering
\caption{Concrete error and risk cases with explanations}
\label{tab:nlp_error_cases}
\begin{tabular}{>{\raggedright\arraybackslash}p{2.9cm} >{\centering\arraybackslash}p{2.0cm} >{\raggedright\arraybackslash}p{2.7cm} >{\raggedright\arraybackslash}p{4.4cm} >{\raggedright\arraybackslash}p{2.6cm}}
\toprule
Text & Expected & Model output & Why the output can be wrong/risky & Mitigation in project \\
\midrule
not bad & Positive & classic: Negative & sparse model over-focuses token bad, weak idiom understanding under negation scope & keep transformer comparison and retain uncertainty band for low-confidence conflict \\
good but late delivery & Negative or issue-flagged & transformer: Positive & contextual model gives too much weight to first positive clause and underweights complaint tail & keep classic negative-first triage for support routing; use issue layer for complaint tags \\
gr8 / thx / idk & Uncertain & both: Uncertain & very short text carries weak sentiment signal and high ambiguity & explicit fallback reason too\_short prevents forced hard label \\
refund\_return (rare label) & Positive detection of rare issue & issue model recall fluctuates & low support label causes unstable boundary and threshold sensitivity & char n-gram variant and per-label threshold tuning; prioritize more annotation on rare labels \\
\bottomrule
\end{tabular}
\end{table}

The main failure patterns are mixed polarity with contrast language, very short low-information text, and rare issue labels in multi-label extraction. The project explicitly tracks these patterns and uses abstention plus per-label controls instead of hiding them behind one aggregate score.

\section{Demo Design for Defense}
The live demo is designed for a 10-minute defense with stable and auditable behavior. The main flow is: input review text, run classic sentiment with uncertainty, optionally run transformer comparison, then trigger issue extraction when review is negative or needs attention.

For UI delivery, the project provides a lightweight interactive app and keeps CLI fallback scripts for reliability. The demo script uses three core scenarios: a clear positive review, a clear negative review with issue tags, and an ambiguous short review that should return uncertain with reason. This scenario-based demo helps examiners verify not only accuracy but also safety behavior and practical usability.

\section{Practical Value for NLP Course}
This project now shows an even wider set of NLP ideas in one system. It covers text normalization, lexical mapping, n-gram sparse encoding, negation scope handling, clause-level cues, multi-label text classification, threshold calibration, and transformer fine-tuning. It now also includes explicit Naive Bayes, Perceptron, and N-gram language model artifacts, plus a feed-forward network baseline on semantic SVD vectors.

The project remains practical for course work. It has saved artifacts, reproducible metrics, an annotation workflow, and deployable inference scripts. It also has a topic-coverage matrix that can be shown directly in defense.

For course learning, the updated version provides stronger evidence of trade-offs between classic and deep NLP approaches under one business objective and one fixed dataset.

\section{Limitations and Future Work}
The main limitation is still data imbalance, especially in negative sentiment and rare issue labels. This affects precision and stability in several models.

Another limitation is topic coverage depth. The project now covers more course topics, but RNN, masked language model tasks, and LLM applications are still light compared with the rest.

Future work can still avoid new data at first. It can add a small RNN baseline on the same split, add a masked language model probe with the same text corpus, and add a simple retrieval or prompt-based LLM application demo. This will improve topic coverage while keeping the same dataset.

\section{Conclusion}
This NLP project builds a complete review understanding pipeline from sentiment classification to issue extraction and transformer extension. Classic context-aware sparse modeling gives strong negative-first performance. Transformer extension gives better global precision and useful semantic behavior in some hard phrases. Multi-label issue extraction provides detailed action signals from review text. The upgraded syllabus block adds Naive Bayes, Perceptron, feed-forward semantic baselines, and N-gram language model evaluation without adding data. With this upgrade, the project gives stronger alignment with course topics and remains reproducible for defense and practical use.

\clearpage
\section*{References}
\addcontentsline{toc}{section}{References}
\begin{thebibliography}{99}

\bibitem{jurafsky2026}
Daniel Jurafsky and James H. Martin.
\textit{Speech and Language Processing} (3rd ed. draft).
Online manuscript, January 2026.
\newline\texttt{https://web.stanford.edu/~jurafsky/slp3/}

\bibitem{manning2008ir}
Christopher D. Manning, Prabhakar Raghavan, and Hinrich Schutze.
\textit{Introduction to Information Retrieval}.
Cambridge University Press, 2008.
\newline\texttt{https://nlp.stanford.edu/IR-book/}

\bibitem{turney2002}
Peter D. Turney.
Thumbs up or thumbs down? Semantic orientation applied to unsupervised classification of reviews.
In \textit{Proceedings of ACL}, 2002.
\newline\texttt{https://aclanthology.org/P02-1053/}

\bibitem{wang2012}
Sida I. Wang and Christopher D. Manning.
Baselines and Bigrams: Simple, Good Sentiment and Topic Classification.
In \textit{Proceedings of ACL}, 2012.
\newline\texttt{https://aclanthology.org/P12-2018/}

\bibitem{wiegand2010}
Michael Wiegand, Alexandra Balahur, Benjamin Roth, Dietrich Klakow, and Andres Montoyo.
A Survey on the Role of Negation in Sentiment Analysis.
In \textit{Proceedings of the Workshop on Negation and Speculation in Natural Language Processing}, 2010.
\newline\texttt{https://aclanthology.org/W10-3111/}

\bibitem{joulin2017}
Armand Joulin, Edouard Grave, Piotr Bojanowski, and Tomas Mikolov.
Bag of Tricks for Efficient Text Classification.
In \textit{Proceedings of EACL}, 2017.
\newline\texttt{https://aclanthology.org/E17-2068/}

\bibitem{nielsen2011}
Finn Arup Nielsen.
A new ANEW: Evaluation of a word list for sentiment analysis in microblogs.
\textit{arXiv preprint arXiv:1103.2903}, 2011.
\newline\texttt{https://arxiv.org/abs/1103.2903}

\bibitem{forman2003}
George Forman.
An extensive empirical study of feature selection metrics for text classification.
\textit{Journal of Machine Learning Research}, 3:1289--1305, 2003.
\newline\texttt{https://www.jmlr.org/papers/v3/forman03a.html}

\bibitem{rifkin2004}
Ryan Rifkin and Aldebaro Klautau.
In defense of one-vs-all classification.
\textit{Journal of Machine Learning Research}, 5:101--141, 2004.
\newline\texttt{https://www.jmlr.org/papers/v5/rifkin04a.html}

\bibitem{tsoumakas2011}
Grigorios Tsoumakas, Ioannis Katakis, and Ioannis Vlahavas.
Random k-Labelsets for multilabel classification.
\textit{IEEE Transactions on Knowledge and Data Engineering}, 23(7):1079--1089, 2011.
\newline\texttt{https://doi.org/10.1109/TKDE.2010.164}

\bibitem{zhang2014}
Min-Ling Zhang and Zhi-Hua Zhou.
A review on multi-label learning algorithms.
\textit{IEEE Transactions on Knowledge and Data Engineering}, 26(8):1819--1837, 2014.
\newline\texttt{https://doi.org/10.1109/TKDE.2013.39}

\bibitem{vaswani2017}
Ashish Vaswani, Noam Shazeer, Niki Parmar, Jakob Uszkoreit, Llion Jones, Aidan N. Gomez, Lukasz Kaiser, and Illia Polosukhin.
Attention Is All You Need.
In \textit{Advances in Neural Information Processing Systems 30}, 2017.
\newline\texttt{https://arxiv.org/abs/1706.03762}

\bibitem{devlin2019}
Jacob Devlin, Ming-Wei Chang, Kenton Lee, and Kristina Toutanova.
BERT: Pre-training of Deep Bidirectional Transformers for Language Understanding.
In \textit{Proceedings of NAACL-HLT}, 2019.
\newline\texttt{https://aclanthology.org/N19-1423/}

\bibitem{saito2015}
Takaya Saito and Marc Rehmsmeier.
The precision-recall plot is more informative than the ROC plot when evaluating binary classifiers on imbalanced datasets.
\textit{PLOS ONE}, 10(3):e0118432, 2015.
\newline\texttt{https://doi.org/10.1371/journal.pone.0118432}

\bibitem{geifman2017}
Yonatan Geifman and Ran El-Yaniv.
Selective Classification for Deep Neural Networks.
\textit{arXiv preprint arXiv:1705.08500}, 2017.
\newline\texttt{https://arxiv.org/abs/1705.08500}

\end{thebibliography}

\clearpage
\section*{Appendix A: Sentiment Figures}
\addcontentsline{toc}{section}{Appendix A: Sentiment Figures}

\begin{figure}[H]
\centering
\includegraphics[width=0.82\textwidth]{\detokenize{results/dm2_steps/01_rating_distribution.png}}
\caption{Rating distribution in source data}
\end{figure}

\begin{figure}[H]
\centering
\includegraphics[width=0.82\textwidth]{\detokenize{results/dm2_steps/06_f1_vs_k.png}}
\caption{Validation F1 by Chi-square top K}
\end{figure}

\begin{figure}[H]
\centering
\includegraphics[width=0.82\textwidth]{\detokenize{results/dm2_steps/06_recall0_vs_k.png}}
\caption{Validation recall for negative class by top K}
\end{figure}

\begin{figure}[H]
\centering
\includegraphics[width=0.82\textwidth]{\detokenize{results/dm2_steps/06b_recall0_by_variant.png}}
\caption{Test recall for negative class by classic NLP variant}
\end{figure}

\begin{figure}[H]
\centering
\includegraphics[width=0.82\textwidth]{\detokenize{results/dm2_steps/06b_precision0_by_variant.png}}
\caption{Test precision for negative class by classic NLP variant}
\end{figure}

\begin{figure}[H]
\centering
\includegraphics[width=0.82\textwidth]{\detokenize{results/dm2_steps/09_prob_hist_test_strong.png}}
\caption{Classic model probability histogram on strong test set}
\end{figure}

\begin{figure}[H]
\centering
\includegraphics[width=0.82\textwidth]{\detokenize{results/dm2_steps/10_threshold_tradeoff.png}}
\caption{Classic model threshold trade-off for coverage and recall}
\end{figure}

\begin{figure}[H]
\centering
\includegraphics[width=0.82\textwidth]{\detokenize{results/nlp_ext/nlp_confusion_matrix.png}}
\caption{Transformer confusion matrix on test set}
\end{figure}

\begin{figure}[H]
\centering
\includegraphics[width=0.82\textwidth]{\detokenize{results/nlp_ext/nlp_prob_hist_test.png}}
\caption{Transformer probability histogram on test set}
\end{figure}

\begin{figure}[H]
\centering
\includegraphics[width=0.82\textwidth]{\detokenize{results/nlp_ext/nlp_threshold_tradeoff.png}}
\caption{Transformer threshold trade-off for coverage and recall}
\end{figure}

\clearpage
\section*{Appendix B: Issue Extraction Figures}
\addcontentsline{toc}{section}{Appendix B: Issue Extraction Figures}

\begin{figure}[H]
\centering
\includegraphics[width=0.82\textwidth]{\detokenize{results/issue_steps/02_label_frequency.png}}
\caption{Issue label frequency in manual annotation set}
\end{figure}

\begin{figure}[H]
\centering
\includegraphics[width=0.82\textwidth]{\detokenize{results/issue_steps/02_per_label_f1.png}}
\caption{Per-label F1 for word issue extraction model}
\end{figure}

\begin{figure}[H]
\centering
\includegraphics[width=0.82\textwidth]{\detokenize{results/issue_steps_char_demo/02_per_label_f1.png}}
\caption{Per-label F1 for word plus char issue extraction model}
\end{figure}

\clearpage
\section*{Appendix C: Syllabus Upgrade Figures}
\addcontentsline{toc}{section}{Appendix C: Syllabus Upgrade Figures}

\begin{figure}[H]
\centering
\includegraphics[width=0.82\textwidth]{\detokenize{results/nlp_ext/syllabus_upgrade/nlp_syllabus_bench_recall0.png}}
\caption{Test recall\_0 across added syllabus baselines}
\end{figure}

\begin{figure}[H]
\centering
\includegraphics[width=0.82\textwidth]{\detokenize{results/nlp_ext/syllabus_upgrade/nlp_syllabus_bench_f2.png}}
\caption{Test F2\_0 across added syllabus baselines}
\end{figure}

\end{document}
